\begin{minipage}[t]{0.48\linewidth}

% === NodeSelector & nodeName ==============================================
\tcbsubtitle{NodeSelector \& nodeName}
\begin{cheatcodebash}
# Add label to node:
k label node <node> disktype=ssd
k label node <node> zone=us-east-1a
\end{cheatcodebash}

\begin{cheatcodebash}
# In pod spec (hard requirement):
nodeSelector:          # <pod>.spec.nodeSelector
  disktype: ssd
\end{cheatcodebash}

\begin{cheatcodebashbold}
# Bypass scheduler entirely in .spec.nodeName:
nodeName: <node-name>  # <pod>.spec.nodeName
\end{cheatcodebashbold}

\begin{cheatcodebash}
### Exam helpers:
k explain pod.spec.nodeSelector
k explain pod.spec.nodeName
# docs: "assign pods nodes"
\end{cheatcodebash}

\end{minipage}\hfill
\begin{minipage}[t]{0.48\linewidth}


% === Static Pods ==========================================================
\tcbsubtitle{Static Pods}
{\footnotesize
  \textbf{kube-apiserver}, \textbf{kube-controller-manager}, \textbf{kube-scheduler}, \textbf{etcd}
  \\
  (\textit{kube-proxy} is a DaemonSet and \textit{CoreDNS} is a Deployment)
}

\vspace{1ex}
\begin{cheatcodebashbold}
# Place manifest in staticPodPath:
/etc/kubernetes/manifests/my-pod.yaml
# Kubelet auto-creates pod; name = <pod>-<node>
# To delete: remove the manifest file
\end{cheatcodebashbold}

\begin{cheatcodebashbold}
# View static pods
k get pods -n kube-system
\end{cheatcodebashbold}


\end{minipage}

\vspace{1.5ex}

\begin{minipage}[t]{0.48\linewidth}

% === Node Affinity =========================================================
\tcbsubtitle{Node Affinity}
\begin{cheatcodebash}
affinity:              # <pod>.spec.affinity
  nodeAffinity:
    # Hard rule (must match):
    requiredDuringSchedulingIgnoredDuringExecution:
      nodeSelectorTerms:
      - matchExpressions:
        - key: disktype
          operator: In     # In|NotIn|Exists|DoesNotExist
          values: [ssd]
    # Soft rule (prefer):
    preferredDuringSchedulingIgnoredDuringExecution:
    - weight: 1
      preference:
        matchExpressions:
        - {key: zone, operator: In, values: [us-east]}
\end{cheatcodebash}

\begin{cheatcodebash}
### Exam helpers:
k explain pod.spec.affinity.nodeAffinity --recursive
# docs: "node affinity"
\end{cheatcodebash}

\end{minipage}\hfill
\begin{minipage}[t]{0.48\linewidth}


% === Pod Affinity / Anti-Affinity =========================================
\tcbsubtitle{Pod Affinity / Anti-Affinity}
\begin{cheatcodebash}
affinity:              # <pod>.spec.affinity
  podAffinity:          # co-locate with matching pods
    requiredDuringSchedulingIgnoredDuringExecution:
    - topologyKey: kubernetes.io/hostname
      labelSelector:
        matchLabels: {app: cache}
  podAntiAffinity:      # spread away from matching pods
    preferredDuringSchedulingIgnoredDuringExecution:
    - weight: 100
      podAffinityTerm:
        topologyKey: kubernetes.io/hostname
        labelSelector:
          matchLabels: {app: web}
\end{cheatcodebash}

\begin{cheatcodebash}
# For each node, scheduler checks every preferred rule. If node satisfies, rule's weight (1-100) adds to node's score. Node with highest total score wins. Required rules have no weight — they're binary.
\end{cheatcodebash}

\begin{cheatcodebash}
### Exam helpers:
k explain pod.spec.affinity.podAffinity --recursive
k explain pod.spec.affinity.podAntiAffinity --recursive
# docs: "inter-pod affinity"
\end{cheatcodebash}

\end{minipage}

\vspace{1.5ex}

\begin{minipage}[t]{0.48\linewidth}


% === Resource Requests & Limits ===========================================
\tcbsubtitle{Resources, LimitRange \& ResourceQuota}
\begin{cheatcodebash}
# CPU units: 1 = 1000m (millicores)
# LimitRange sets namespace defaults:
apiVersion: v1
kind: LimitRange
metadata: {name: lr, namespace: <ns>}
spec:
  limits:
  - type: Container
    # limit applied none specified (resources.limits)
    default:        {cpu: 500m, memory: 256Mi}
    # request applied none requested (resources.requests)
    defaultRequest: {cpu: 100m, memory: 128Mi}
    # maximum allowed limit a container may set
    max:            {cpu: "2",  memory: 1Gi}
\end{cheatcodebash}
\begin{cheatcodebash}
# ResourceQuota caps total namespace usage:
apiVersion: v1
kind: ResourceQuota
metadata: {name: rq, namespace: <ns>}
spec:
  hard:
    pods: "10"
    requests.cpu: "4"
    requests.memory: 8Gi
    limits.cpu: "8"
    persistentvolumeclaims: "5"
\end{cheatcodebash}

\end{minipage}\hfill
\begin{minipage}[t]{0.48\linewidth}

% === Taints & Tolerations =================================================
\tcbsubtitle{Taints \& Tolerations}
\begin{cheatcodebash}
# Taint effects:
# NoSchedule        - no new pods w/o toleration
# PreferNoSchedule  - soft NoSchedule
# NoExecute         - evict existing pods too
k taint node <node> key=val:NoSchedule
k taint node <node> key=val:NoExecute
k taint node <node> key-       # remove
# Node has this taint:                                                                                                                                                                                
k taint node <node> dedicated=gpu:NoSchedule
\end{cheatcodebash}

\begin{cheatcodebash}
  # Pod must have this toleration to be scheduled there:
  tolerations:           # <pod>.spec.tolerations[]
  - key: "dedicated"
    operator: "Equal"    # Exists
    value: "gpu"
    effect: "NoSchedule"
\end{cheatcodebash}

\note{\footnotesize
  \begin{itemize}
    \item With \textbf{affinity} one selects nodes by their labels
    \item With \textbf{taints/tolerations} the node repels pods, and the pod explicitly opts in by declaring a matching toleration.
  \end{itemize}
}

\begin{cheatcodebash}
### Exam helpers:
k explain node.spec.taints --recursive
k explain pod.spec.tolerations --recursive
# docs: "taints tolerations"
\end{cheatcodebash}

\end{minipage}
